\documentclass[letterpaper,12pt]{article}
%\usepackage{harvard}
%\usepackage{subfigure}
%\usepackage[active]{srcltx}
\usepackage{graphicx}
\usepackage{pdflscape}
\usepackage{amsmath}
\usepackage{amsthm}
\usepackage{lscape}
\usepackage{enumerate}
\usepackage{float}
\usepackage{grffile}
\usepackage{enumitem}
\usepackage{relsize}
\usepackage{bm}
\usepackage{tablefootnote}
\usepackage{footnote}
\usepackage{sectsty}
\usepackage{xcolor}
\usepackage{caption}
\usepackage{color}
\usepackage{natbib}
\usepackage{multirow}
\usepackage{bbm}
\usepackage[T1]{fontenc}
\usepackage[multiple]{footmisc}
\usepackage{appendix}
\usepackage{times}
\usepackage{relsize}
\usepackage{caption}
\usepackage{newtxtext,newtxmath}
\usepackage{enumitem}


%%%%%%%%%%%%%%%%%%%%%%%%%%%%%%%%%%%%%%%%%%%%%%%%%%%%%%%%%%%%
%%%   Title Page Information
%%%%%%%%%%%%%%%%%%%%%%%%%%%%%%%%%%%%%%%%%%%%%%%%%%%%%%%%%%%%
\title{\vspace{-1.2in} ECON712: Final Exam}

\date{\vspace{-0.5 in} April 16th, 2026}


%%%%%%%%%%%%%%%%%%%%%%%%%%%%%%%%%%%%%%%%%%%%%%%%%%%%%%%%%%%%
%%%   Double and Single Space Commands
%%%%%%%%%%%%%%%%%%%%%%%%%%%%%%%%%%%%%%%%%%%%%%%%%%%%%%%%%%%%
\newlength{\defbaselineskip}
\setlength{\defbaselineskip}{\baselineskip}
\newcommand{\setlinespacing}[1]{\setlength{\baselineskip}{#1 \defbaselineskip}}
\newcommand{\doublespacing}{\setlength{\baselineskip}{1.3 \defbaselineskip}}
\newcommand{\singlespacing}{\setlength{\baselineskip}{1\defbaselineskip}}
\renewcommand{\thesubsection}{\arabic{subsection}}
\newtheoremstyle{dotlessS}{}{}{\itshape}{}{\color{dg}\bfseries}{}{ }{}
\theoremstyle{dotlessS}
\newtheorem{theorem}{Theorem}
\newtheorem{lemma}{Lemma}
\newtheorem{proposition}{Proposition}
\newtheorem{assumption}[theorem]{Assumption}
\newtheorem{corollary}{Corollary}
\newtheorem{definition}{Definition}
\newenvironment{example}[1][Example]{\begin{trivlist}
\item[\hskip \labelsep {\bfseries #1}]}{\end{trivlist}}
\newenvironment{remark}[1][Remark]{\begin{trivlist}
\item[\hskip \labelsep {\bfseries #1}]}{\end{trivlist}}

\newcommand\Pitem{%
  \addtocounter{enumi}{-1}%
  \renewcommand\theenumi{\arabic{\enumi}'}%
  \item%
  \renewcommand\theenumi{\arabic{enumi}}%
}

%%%%%%%%%%%%%%%%%%%%%%%%%%%%%%%%%%%%%%%%%%%%%%%%%%%%%%%%%%%%
%%%   Set Margins Here
%%%%%%%%%%%%%%%%%%%%%%%%%%%%%%%%%%%%%%%%%%%%%%%%%%%%%%%%%
\textwidth=6.5in
\textheight=9in
\oddsidemargin=0.10in
\evensidemargin=0.10in
\topmargin=-0.5in
%\parskip=3pt

%%%%%%%%%%%%%%%%%%%%%%%%%%%%%%%%%%%%%%%%%%%%%%%%%%%%%%%%%%%%%
%%%   Actual Document Begins Here
%%%%%%%%%%%%%%%%%%%%%%%%%%%%%%%%%%%%%%%%%%%%%%%%%%%%%%%%%%%%%
\begin{document}
\pagenumbering{arabic}
\maketitle
\vspace{-.2in}
\noindent
\vspace{-0.6 in}

\section*{ Question I: Heterogeneity and Causal Inference}
Suppose that, for a sample of individuals, average cancer rates of smokers and non-smokers are given as follows:
\begin{table}[H]
\begin{center}
\caption{Cancer Diagnosis by Age Group and Smoking Status}
\label{table:main_shock_effect}
\resizebox{.42\textwidth}{!}{%
\begin{tabular}{lll}
\hline
\hline
Age Group (x)   &  Smoker           & Non-Smoker  \\
\hline
60-70  & 75 (80)                      & 25 (40) \\
50-60 &  31 (40)                      &10 (20)  \\
40-60  & 6 (10)                 & 10 (90)  \\
30-40 & 2 (5)                 &   3 (45)  \\
\hline
\hline
\end{tabular}}
\end{center}
\captionsetup{font={scriptsize}}
\captionsetup{width=0.80\textwidth}
\captionsetup{skip=0pt}
\caption*{Notes: Sample size are given in brackets}
\end{table}
\noindent Suppose that you want to estimate the following regression:
\begin{eqnarray*}
y_i=\beta_0+\beta_1 \: D_i+\epsilon
\end{eqnarray*}
where:
\begin{center}
\begin{minipage}{0.45\textwidth}
\[
D_i=\begin{cases}
		1 & \text{if person $i$ is a smoker} \\
             0 & \text{if person $i$ is a non-smoker}
\end{cases}
\]
\end{minipage}
\hfill
\begin{minipage}{0.45\textwidth}
\[
y_i=\begin{cases}
		1 & \text{if person $i$ has cancer} \\
             0 & \text{if person $i$ does not have cancer}
\end{cases}
\]
\end{minipage}
\end{center}
Using the information in this table, answer the following questions:
\begin{enumerate}[label=(\alph*)]
\item (7 points) Suppose that you use the sample to estimate the regression, calculate $\hat{\beta}_1^{OLS/MM}$. Is this an unbiased estimator of ATE? Is this a consistent estimator of ATE? Explain.
\item (7 points) \textit{Assume that $y(0),y(1) \perp D|x$}, calculate $ATE$ and $ATT$
\end{enumerate}

\section*{Question II: OLS Finite Sample Properties
}
Assume that the model $\bm{y}=\bm{X} \: \bm{\beta}+\bm{u}$ satisfies the Gauss-Markov assumptions (i.e. classical assumptions), where $\bm{X}$ is a $n x k$ matrix of observations. Define a new matrix $\bm{Z}=a \: \bm{X}$, where $a>1$ is a scalar. Suppose a researcher draws a random sample $n$ and obtains the following estimator:
\begin{eqnarray*}
\bm{\widetilde{\beta}}=(\bm{Z}'\bm{X})^{-1}\bm{Z}'\bm{y}
\end{eqnarray*}
Using this information answer the following question:
\begin{enumerate}[label=(\alph*)]  
\vspace{0.11in}
\item (5 points) Prove whether $\bm{\widetilde{\beta}}$ is a biased or unbiased estimator of $\bm{\beta}$? 
\vspace{0.11in}
\item (5 points) Calculate the variance-covariance matrix of the estimator. Which estimator do you prefer, $\bm{\widetilde{\beta}}$ or $\bm{\hat{\beta}}$, where $\bm{\hat{\beta}}=(\bm{X}'\bm{X})^{-1}\bm{X}'\bm{y}$ is the method-of-moment estimator derived in class? Explain.
\vspace{0.11in}
\item (6 points) Suppose now that you use the OLS estimator, $\bm{\hat{\beta}}$, but that there is heteroscedasticity. In particular, $E[\bm{u}\bm{u}']=\Omega$, where:
\begin{eqnarray*}
\bm{\Omega}=
 \begin{bmatrix}
    \sigma^2_1 & 0  & \cdots &0        \\
     0           & \sigma^2_2  & \cdots &0 \\
     \vdots  &\vdots           &\vdots &\vdots \\ 
      0           & 0  & \cdots &  \sigma^2_N  \\
\end{bmatrix}
\end{eqnarray*}
Carefully explain how you would estimate the variance-covariance matrix of the OLS estimator when $\Omega$ is not known (make sure to discuss the importance of the sandwich estimator in your answer). Is OLS the most efficient estimator in this case?   
\end{enumerate}


\section*{Question III: Panel Data and Difference-in-Differences}

Consider a balanced panel dataset with $n$ units observed over $T=2$ periods ($t \in \{1,2\}$). The true causal population regression is:
\begin{equation*}
    y_{it} = \beta x_{it} + \alpha_i + \lambda_t + \varepsilon_{it}, \qquad i=1,\ldots,n,\quad t \in \{1,2\},
\end{equation*}
where $x_{it}$ is a scalar regressor, $\alpha_i$ is a \textbf{unit fixed effect} that may be correlated with $x_{it}$, $\lambda_t$ is a \textbf{common time effect} shared by all units (e.g.\ a macroeconomic shock), and $\varepsilon_{it}$ is an idiosyncratic error. Assume strict exogeneity throughout: $\mathbb{E}[\varepsilon_{it} \mid x_{i1}, x_{i2}, \alpha_i] = 0$ for $t \in \{1,2\}$. In addition, note that since $\lambda_t$ is a fixed parameter, the conditional mean of the outcome net of the unit effect and the regressor is period-specific:
\begin{equation*}
    \mathbb{E}\!\left[y_{it} - \beta x_{it} - \alpha_i \;\middle|\; \bm{x}_i,\, t=1\right] = \lambda_1, \qquad \mathbb{E}\!\left[y_{it} - \beta x_{it} - \alpha_i \;\middle|\; \bm{x}_i,\, t=2\right] = \lambda_2,
\end{equation*}
where $\lambda_1$ and $\lambda_2$ are not restricted to be equal, capturing a common time effect that may differ across periods.

\medskip
\noindent Stack all $2n$ observations into the $(2n \times 1)$ vectors $\bm{y}$, $\bm{x}$, $\bm{\varepsilon}$, and $\bm{\lambda}$ (where $\lambda_t$ is repeated $n$ times for each period), and let $\bm{D}$ be the $(2n \times n)$ matrix of unit dummies. Define the composite error $\bm{u} \equiv \bm{D}\bm{\alpha} + \bm{\lambda} + \bm{\varepsilon}$. The true model in stacked form is:
\begin{equation*}
    \bm{y} = \bm{x}\beta + \bm{D}\bm{\alpha} + \bm{\lambda} + \bm{\varepsilon}.
\end{equation*}

\begin{enumerate}[label=(\alph*)]

\item (7 points) A researcher who ignores both $\alpha_i$ and $\lambda_t$ runs pooled OLS of $\bm{y}$ on $\bm{x}$ and a constant term, treating $\bm{u}$ as the error term. Derive the pooled OLS estimator $\hat{\beta}^{OLS}$ in matrix form as a function of the true population parameter $\beta$ and fixed effect terms. Carefully discuss why the pooled OLS might be inconsistent \textbf{(you do not need to show the math)}.
\item (7 points) The researcher now attempts to correct for unit fixed effects by applying the within transformation. For each unit $i$, define demeaned variables:
\begin{equation*}
    \tilde{y}_{it} = y_{it} - \bar{y}_i, \qquad \tilde{x}_{it} = x_{it} - \bar{x}_i,
\end{equation*}
where $\bar{y}_i = \frac{1}{2}(y_{i1}+y_{i2})$ and $\bar{x}_i = \frac{1}{2}(x_{i1}+x_{i2})$. By the Frisch--Waugh--Lovell theorem, running OLS on the demeaned variables $\tilde{\bm{y}}$ and $\tilde{\bm{x}}$ is equivalent to partialling out unit fixed effects from both $\bm{y}$ and $\bm{x}$. Write the demeaned model for $\tilde{y}_{it}$ in terms of $\tilde{x}_{it}$, $\tilde{\varepsilon}_{it}$, and the demeaned time effects $\tilde{\lambda}_t \equiv \lambda_t - \bar{\lambda}$, where $\bar{\lambda}=\frac{1}{2}(\lambda_1+\lambda_2)$. Show that the common time effect does \emph{not} vanish after the within transformation and express $\tilde{\lambda}_1$ and $\tilde{\lambda}_2$ in terms of $\delta \equiv \lambda_2 - \lambda_1$.  Provide an intuitive explanation of why controlling for unit fixed effects might not sufficient.



\item (7 points) Suppose now that individuals $i$ belong to one of two groups $s \in \{0,1\}$, a control group ($s=0$) and a treatment group ($s=1$), observed over $t \in \{1,2\}$. Let $D_s$ be a dummy variable equal to one if the individual belongs to the treatment group and zero otherwise, and let $\lambda_t$ be a dummy variable equal to one if the observation is in the second period and zero otherwise. Consider the following population regression:
\begin{equation*}
    y_{ist} = \alpha + \beta_1 D_s + \beta_2 \lambda_t + \beta_3 \left(D_s \cdot \lambda_t\right) + u_{ist}.
\end{equation*}
Discuss the difference-in-differences estimator of $\beta_3$, the assumptions required for $\beta_3$ to have a causal interpretation, and express $\hat{\beta}_3$ in terms of group-period sample means $\bar{y}_{s,t}$.

\end{enumerate}


\section*{Question IV: Linear transformation of variables (8 points)}
Assume the regression model:
\begin{eqnarray*}
\bm{y}=\bm{\iota_n} \: \hat{\gamma} + \bm{X} \: \bm{\hat{\beta}}+\bm{\hat{u}}
\end{eqnarray*}
where
\begin{itemize}
\item $\bm{y}$ is a $nx1$ vector of dependent variables
\item $\bm{\iota_n}$ is an $nx1$ vector of ones, 
\item $\bm{X}$ is a $nxk$ matrix of regressors, 
\item $\hat{\gamma}$ is a scalar OLS/MM estimate
\item $\bm{\hat{\beta}}$ is a $kx1$ vector of OLS/MM estimates
\item $\bm{\hat{u}}$ is a $nx1$ vector of residuals
\end{itemize}
Define the following transformed variables:
\begin{eqnarray*}
\bm{z}=\bm{y}- \bm{\iota_n} \overline{y}
\end{eqnarray*}
and
\begin{eqnarray*}
\bm{W}=\bm{X}-\bm{\iota_n} \bm{\overline{x}}^T
\end{eqnarray*}
with
\begin{eqnarray*}
\begin{array}{lllll}
\overline{y}=\frac{1}{n}\underset{i=1}{\overset{n}{\sum}}y_{i} &,& \bm{\overline{x}}=\begin{pmatrix} \bar{x}_1 \\ \bar{x}_2 \\ \vdots  \\  \bar{x}_k \end{pmatrix} &,& \bar{x}_j=\frac{1}{n}\underset{i=1}{\overset{n}{\sum}}x_{ij} 
\end{array}
\end{eqnarray*}
Suppose that a researcher uses OLS/MM to estimate the following regression isntead:
\begin{eqnarray*}
\bm{z}= \bm{W} \: \bm{\hat{\alpha}}+\bm{\hat{\epsilon}}
\end{eqnarray*}
Derive the relationship between $\bm{\hat{\alpha}}$ and $\bm{\hat{\beta}}$. Are the residuals from these two regressions equal?





\section*{Question V: Local Average Treatment Effect}

Consider a setting with a binary outcome $Y_i$, a binary treatment $D_i \in \{0,1\}$, and a binary instrument $Z_i \in \{0,1\}$. Let $Y_i(d)$ denote the potential outcome under treatment status $d$, and let $D_i(z)$ denote the treatment status individual $i$ would take under instrument value $z$. The Wald estimator is defined as:
\begin{equation*}
    \hat{\beta}^{IV} = \frac{\mathbb{E}[Y_i \mid Z_i = 1] - \mathbb{E}[Y_i \mid Z_i = 0]}{\mathbb{E}[D_i \mid Z_i = 1] - \mathbb{E}[D_i \mid Z_i = 0]}.
\end{equation*}
Assume the following four conditions hold:
\begin{enumerate}[label=(\roman*)]
    \item \textbf{Independence:} $(Y_i(0), Y_i(1), D_i(0), D_i(1)) \perp Z_i$
    \item \textbf{Exclusion restriction:} $Z_i$ affects $Y_i$ only through $D_i$
    \item \textbf{Relevance:} $\mathbb{E}[D_i \mid Z_i=1] \neq \mathbb{E}[D_i \mid Z_i=0]$
    \item \textbf{Monotonicity:} $D_i(1) \geq D_i(0)$ for all $i$
\end{enumerate}

\begin{enumerate}[label=(\alph*)]

\item (5 points) Under the assumptions above, prove that:
\begin{equation*}
    \hat{\beta}^{IV} = \mathbb{E}\!\left[Y_i(1) - Y_i(0) \;\middle|\; D_i(1) > D_i(0)\right].
\end{equation*}
Carefully indicate the assumption that allows you to write the equation differently at each stage of the proof.

\item (6 points) Discuss the external validity of $\hat{\beta}^{IV}$. For which subpopulation does LATE recover the treatment effect, and under what conditions does it coincide with the Average Treatment Effect? What are the implications for policy evaluation?

\end{enumerate}


\end{document}
